\section{Discussion}\label{sec:discussion}

\subsection{Cross-tier comparison}
\label{sec:disc-cross-tier}

The core inverse-MC scoring machinery works cleanly across all three
tiers: the inverter, FD-checked gradient, contrast diagnostic, and
feature plumbing operate unchanged on a synthetic $n=100$ chain, a SUMO
$n=2{,}405$ network, and a real-world $n=1{,}733$ Xuancheng network. What
differs by tier is the loader, the regime definition, and the
OD-control logic; the shared scoring core is selected by a single
\code{--dataset \{sumo,xuancheng\}} command-line interface (CLI) flag. This is a useful
generalisation result for an inverse-MC pipeline that the original paper
validated only on synthetic data and one real-world demonstration.

What does not generalise as cleanly is the magnitude of the
detection signal. Tier~1 reproduces the paper's RMAE within
$\sim 0.05$ for $|\Xo| \ge 10$. Tier~2 reaches $\AUC = 0.708$ at $T=40$
on the controlled SUMO comparison, $\sim 69\%$ of the gap to the
full-observation empirical $1$-step reference of $0.801$, with per-row
chain Pearson $\sim +0.21$. Tier~3 on Xuancheng, with no ground-truth
sat-navigation label, gives chance-level fitted AUC on the
rush-vs-off-peak behavioural proxy. A stronger Labour-Day proxy creates
a significant fine-OD empirical signal ($\AUC_{\mathrm{emp}}=0.567$),
but multistart validation-selected fitting still falls to chance
($\AUC_{\mathrm{fit}}=0.502$). The drop
between Tier~2 and Tier~3 has two distinct contributions:
\begin{itemize}
\item The Xuancheng ceiling itself ($\AUC_{\mathrm{emp}} \sim 0.57$) is
much lower than the SUMO ceiling ($\sim 0.69$), reflecting that the
behavioural rush-vs-off-peak proxy is a weaker contrast than the
SUMO labelled comparison.
\item Xuancheng is also scored on shorter legalised trajectories
($T=10$), whereas the main SUMO detector starts at $T=20$ and improves
again at $T=30$ and $T=40$. Since the LLR score accumulates evidence over
the trajectory, the shorter real-data window likely suppresses both the
empirical reference and fitted AUC.
\item The fitted-vs-reference gap is not merely a scaled-down
version of the SUMO success. The pooled un-matched Xuancheng row gives
a superficially moderate gap-closure number because the reference itself
is weak, but once OD mix is controlled the fitted detector sits at chance
across every tested $K$ while the full-observation reference remains
weakly above chance at finer granularities. Tier~3 is therefore a
negative result for the current $1$-step partial-observation filter
under this behavioural proxy.
\item The Labour-Day stress test rules out the strongest alternative
reading: that a more intuitive real-world contrast would rescue the
detector. Labour Day does yield a stronger empirical reference after
fine OD matching, but the fitted chain remains near chance and multistart
selection removes the small zero-init lift.
\item A fitting-free $2$-step backoff audit on the same Labour-Day split
also fails to lift the empirical reference: the best setting is the
degenerate $\lambda=0$ backoff, equivalent to the original $1$-step
chain. The SUMO route-memory result therefore motivates a future model
class, but it is not by itself an explanation of the Xuancheng failure.
\end{itemize}
Under controlled labelled simulation the method recovers a robust
fraction of the per-row chain contrast. On a real-world behavioural
proxy, the fitted $1$-step detector remains near chance once
population-mix is controlled.

\Cref{tab:cross-tier} consolidates the three tiers on a common axis:
the same inverter and diagnostic produce every row, and the trend in the
empirical ceiling, from a noise-free synthetic recovery, through a
labelled simulated comparison, to a behavioural real-world proxy,
quantifies how much detectable signal each setting actually contains.

\begin{table}[H]
\centering
\caption{Cross-tier synthesis. Tier~2 reports the representative
Ingolstadt seed~$42$ at $T=20$; Xuancheng rows report the weekday proxy
and the Labour-Day stress test. Fitted and Pearson cells are
\code{none}/\code{real}.}
\label{tab:cross-tier}
\scriptsize
\setlength{\tabcolsep}{3pt}
\begin{tabular}{@{}p{0.17\linewidth}p{0.23\linewidth}ccc@{}}
\toprule
tier & dataset ($n$) & empirical ceiling & fit AUC & Pearson \\
\midrule
1 synthetic              & random graph ($100$)  & --     & --             & -- \\
2 SUMO                   & Ingolstadt network ($2{,}405$) & $0.685$ & $0.612 / 0.644$ & $+0.132 / +0.215$ \\
3 Xuancheng              & AVI ($1{,}733$)        & $0.573$ & $0.524 / 0.536$ & $+0.096 / +0.072$ \\
3 Xuancheng Labour       & AVI ($1{,}733$)        & $0.567$ & $0.521 \rightarrow 0.502$ & $+0.066$ \\
\bottomrule
\end{tabular}
\end{table}

\subsection{Limits of the chosen approach}
\label{sec:disc-limits}

The main limitations are label availability, stationarity, and model
capacity.

\paragraph{Proxy labels and stationarity.} Tier~3 cannot observe
sat-navigation use directly, so rush-vs-off-peak and Labour-Day splits
are behavioural proxies. They mix route guidance, ordinary congestion
adaptation, driver composition and OD mix. The estimator also compresses
dynamic time windows into one stationary surrogate chain, which is useful
for scoring but should not be interpreted as a literal invariant traffic
distribution.

\paragraph{Controlled simulation versus real traffic.} Tier~2 provides a
labelled test bed, yet it still contains simulator-specific structure:
at the selected demand the sat-navigation regime completes many more
trips, and unaided SUMO drivers follow precomputed routes. The
$\rho \times d$ audit in Section~\ref{sec:phase4-rho-demand} in the
Appendix shows that this simulator signal is mainly tied to the
sat-navigation label, but a real experienced driver can adapt without
live route guidance. This is why Tier~3 remains a proxy stress test until
trip-level guidance labels are available.

\paragraph{Sparse observation, optimisation and feature capacity.} Both
SUMO and Xuancheng show that the inverse problem is under-constrained at
city scale. Multi-seed SUMO runs make the feature result weaker as an AUC
claim but robust as a per-row chain-recovery effect; Xuancheng multistart
removes an apparent \code{features=real} ceiling violation and returns the
fitted detector to chance. The $f+g$ extension adds another capacity
route, but its city-scale implementation remains expensive because each
iteration needs large linear solves. The current thesis therefore uses
the tractable $f$-only detector and treats $f+g$ and higher-order dynamics
as model-class extensions.

\subsection{Social, environmental, technological implications}
\label{sec:disc-implications}

\paragraph{Transport-policy modelling.} Navigation platforms can change
the routes that appear in revealed-preference data. If those influenced
routes are fed back into demand models as ordinary preferences, policy
evaluation can confuse platform response with driver preference. The
controlled SUMO result shows that inverse-MC filtering can recover part
of this distinction when labels exist; the Xuancheng result shows that
behavioural proxies alone are too weak for confident attribution.

\paragraph{Privacy and data access.} The evidence needed for this task is
awkward: trip-level trajectories are useful, but trip-level app-use labels
are sensitive. Aggregate-only releases protect privacy but cannot support
the detector tested here; joinable per-trip releases enable the method
while increasing reidentification risk. A practical deployment would need
privacy-preserving access to coarse guidance/compliance labels, or a
trusted-analysis setting in which labels can be used without publishing
raw identifiers.

\paragraph{Platform accountability.} A reliable detector would let cities
audit how much of observed route choice is shaped by guidance systems,
without needing to inspect individual recommendations. This matters for
congestion pricing, traffic-calming decisions and claims about network
welfare. The negative Tier~3 result is part of that accountability story:
it warns against treating rush-hour AVI contrasts as evidence of
sat-navigation influence unless the proxy is supported by labels or a
stronger behavioural model.

\subsection{Future work}
\label{sec:disc-future}

\begin{enumerate}
\item \textbf{Labelled real-world validation.} The highest-value next
dataset would pair trajectories with coarse sat-navigation-use or
route-guidance-compliance labels. Even partial labels would convert
Tier~3 from a proxy stress test into a direct validation exercise.

\item \textbf{Robustness over stations, OD cells and cities.} The current
fits condition on one observation-station set and a fixed OD match. A
deployment study should repeat the inverse-MC fits over multiple $\Xo$
draws, resample OD cells, and replicate across more city/month pools.

\item \textbf{Richer behavioural state.} SUMO shows that second-order
route memory can carry signal, while Xuancheng Labour does not show a
simple empirical $2$-step lift. A parametric extension should therefore
test higher-order context, time-of-day conditioning and trip-level
features as competing explanations, with the chain-recovery claim kept
separate from any deployment-oriented classifier.

\item \textbf{Scalable $f+g$ optimisation.} The hitting-rate extension
remains attractive because it uses more of Morimura's original evidence,
but city-scale fitting needs factor sharing or a different optimiser.
Sherman--Morrison updates for the $(I-\beta P_T)$ solves and
natural-gradient or Levenberg--Marquardt steps are the most direct
algorithmic candidates.

\item \textbf{Experienced-driver simulation.} The Tier~2
$\rho \times d$ decomposition should be repeated with unaided drivers who
avoid known rush-hour bottlenecks. That counterfactual would make the
simulator bound closer to the real Xuancheng ambiguity between human
experience and algorithmic guidance.
\end{enumerate}
